% The ConceptBase User Manual (CB-Manual)
% Copyright 1987-2023 by The ConceptBase Team
% The CB-Manual may only be updated by members of the ConceptBase Team
% as listed on http://conceptbase.cc
% The TeX sources of the CB-Manual and the accompanying
% figures can be freely distributed under the provision that this
% clause is included in file CB-Manual.tex and that CB-Manual.tex
% is the main file of the TeX sources.
% Contact the ConceptBase Team via the URL above if you want to join
% the team.

% Creation:    7-Mar-2002 Manfred Jeusfeld (UNITILB)
% Last Change: %G% , Manfred Jeusfeld (UNITILB)
% Release:     %R%
%
% -----------------------------------------

\chapter{The ConceptBase.cc Server}
\label{cha:cbserver}


The ConceptBase.cc server (CBserver) offers its services via a TCP/IP port to client
programs.
The main services are to TELL or UNTELL O-Telos objects
and to ASK queries to the database. The operations are called by clients
(for example, the user interfaces described in section \ref{cha:workbench}).
An arbitrary number of clients can connect to a CBserver.

The CBserver is started%
\footnote{You can also start the CBserver from within the ConceptBase.cc user interface CBIva. Details are in the 
installation guide distributed with ConceptBase.cc and in section \ref{sec:session}.}
 by a command line

\begin{verbatim}
   cbserver <params>
\end{verbatim}

assuming that the installation directory of ConceptBase.cc is
added to the search path of executable programs. If it is not on the
search path, then simply change the current directory to the installation
directory of ConceptBase.cc or use its absolute path of the cbserver script.


\section{CBserver parameters}
\label{sec:cbsparams}

The following parameters are available for the 'cbserver' command:

\begin{description}
\item[-d dbdir] Set database {\tt dbdir} to be loaded. If the database does not exist, it
is created and initialized with the O-Telos pre-defined objects.
The database is maintained as a directory. Setting the database is mandatory
except when the update mode is set to {\tt nonpersistent} (see below).
You cannot start two concurrent servers which use the same database directory.
To avoid this case, a file {\tt OB.lock} is created in the database directory
when the first server is started. If the server crashes during its execution,
the file {\tt OB.lock} will still exist in the directory. Before you
restart the server, you might have to remove this file manually.


\item[-db dbdir] Like {\tt -d} but also sets the load/save/views directories to
{\tt dbdir}, i.e. the CBserver will automatically maintain the module sources
in {\tt dbdir} and also materialize the selected queries in the same directory.
See section \ref{sec:module_views} for details.

\item[-new dbdir] Like {\tt -d} but deletes any existing database at location
{\tt dbdir} before it creates and initializes it.

\item[-u updatemode] controls update persistency.
The allowed values are {\tt persistent}  and {\tt non\-per\-si\-stent}.
If no database is provided by parameter "-d", then the default
update mode is set to nonpersistent. Otherwise, the default is
persistent.
In nonpersistent mode, all updates are lost after the ConceptBase server
is stopped. In persistent mode, updates are stored in the files of the database and
will be available for future sessions.

\item[-U untellmode] controls the way how UNTELL is executed by the server.
The allowed values are {\tt verba\-tim} and {\tt cleanup} (default).
In verbatim mode, the UNTELL operation will only untell the facts 
directly described by the O-Telos frame being submitted as argument.
In cleanup mode, UNTELL will also try to remove the instantiation to the
O-Telos system classes {\tt Individual}, {\tt Attribute}, {\tt InstanceOf} and {\tt IsA}. 
By doing so, UNTELL behaves inverse to the TELL
operation. More details are explained in subsection \ref{sec:untell}.

\item[-port portnr] sets the TCP/IP socket portnumber for client connections
to the CBserver. The value {\tt portnr}
must be between 2000 and 65535. If there is already a process using the
portnumber, the CBserver will abort. The default value for
the portnumber is 4001.

\item[-p portnr] is the same as "-port portnr". Deprecated since it conflicts with
a predefined command line parameter of SWI-Prolog.


\item[-version] display version info and exit.

\item[-help] display list of CBserver options and exit.

\item[-license] display license and exit.

\item[-t tracemode] sets the tracemode of the CBserver.
It is one of {\tt silent}, {\tt no}, {\tt minimal}, {\tt low}, {\tt high}, {\tt veryhigh}. The tracemode
determines the amount of text displayed by the server during its execution.
The tracemode does not influence the function but is used for debugging. The
default tracemode is set to {\tt no} (only display CBserver interface).
The tracemode {\tt low} will configure the CBserver to trace the CBserver interface
calls plus answers, and the tracemode {\tt no} virtually disables tracing.
The tracemode {\tt silent} is even surpressing the message 'CBserver ready' when starting up
the CBserver.
The tracemode {\tt high} and {\tt veryhigh} are useful for debugging the system.
In these two modes, an unlikely fatal signal like division by zero will not directly abort the
CBserver process but start a debug dialog. Enter "h" for options to
diagnose the problem in collaboration with the ConceptBase developers.

\item[-c cachemode] turns on the query cache to allow recursive query evaluation.
The value {\tt cachemode} is one of {\tt off}, {\tt transient}, and {\tt keep} (default).
In transient mode the cache is emptied before each transaction. In keep mode,
the cache is emptied when the maximum number of entries in the cache is exceeded
or an update has invalidated the cache. Further details are explained in section \ref{sec:cbcache}.

\item[-cs size] specifies maximum number of derived facts retained in a cache between two transactions.
This option may be used in conjunction with the cachemode  {\tt keep}. The default value is
60000 facts.

\item[-o optmode] controls the optimizer for rules, constraints and queries.
The value {\tt optmode} is one of {\tt 0} (no optimization), {\tt 1}
(structural optimization by exploiting builtin O-Telos axioms),
{\tt 2} (optimizing join order), {\tt 3} (combines 1 and 2), or
{\tt 4} (combines 1 and 2 with trigger pruning). Default
and recommended is 4.

\item[-r secs] automatically restarts the CBserver after a crash, or when it was
started with option -sm slave and the last client exits.
The value {\tt secs} specifies how many seconds to wait before restart.
You may want to use this option in a multi-user setting, where the CBserver
runs on a different machine that the user clients. The -r option must be handled with
great care since it can easily lead to an infinite loop of restarts, e.g. when
a database file is corrupted. In such cases you might have to reboot the whole computer!


\item[-s securitylevel] configures the access control mechanism of
ConceptBase. The value {\tt 0} means that no access control is employed.
Any user can ask, tell, untell, retell any object in any module. You
can also untell objects defined in a super-module.
Level {\tt 1} (default)  provides a very basic protection: one can only
untell objects if they are defined in the current module. This
prevents in particular undesired deletions of objects defined in
the {\tt System} module.
Level {\tt 2} fully enables access control. First, untelling of objects must
happen in the module where the object has been defined. Second, any
transaction submitted by a user to the CBserver is checked against the
permission rules as defined in section \ref{sec:modaccess}.
Level {\tt 3} enables at most read access to a module. In addition, the permission rules
must allow read access. This level makes sense if you want to freeze a database state.
Enable access control
when you use ConceptBase in a multi-user setting and you want to avoid errorneous
interferences between different users.

\item[-e maxerrors] sets the maximum number of errors to be displayed to a ConceptBase.cc
client within one transaction. The value -1 means that no restriction is applied.
Set to 0 to surpress any errors message and to a positive number to limit the
displayed errors messages to that number. A low positive number can speed up the
communication between ConceptBase client and server if a lot of error messages
are generated. The default is 20.

\item[-cc ccmode] (predicate typing) controls to which extent the CBserver
applies strict typing of
attribution predicates {\tt (x m y)} occurring in the membership constraints
of query classes.
If the mode is set to {\tt strict} (=default), attribution predicates without a unique
concerned class%
\footnote{The concerned class is a consequence of the predicate typing condition
of section \ref{sec:CBL}. You can roughly compare it to typing of variables
in programming languages.}
shall not be accepted.
If the mode is set to {\tt extended}, the search for concerned classes shall include
subclasses (see section \ref{sec:SemFormula}). 
If the mode is set to {\tt off}, ConceptBase.cc
also accepts queries with unstrictly-typed attribution predicates.
The strict mode is preferable since it avoids certain semantic errors.
Deduction rules and integrity constraints may never violate the
predicate typing condition, even if the mode is set to {\tt off}.
An example for a query using non-strict predicate typing is available from 
the CB-Forum, see \url{http://merkur.informatik.rwth-aachen.de/pub/bscw.cgi/1270138}.

\item[-mu mumode] (multi-user mode) specifies whether the CBerver should
run in multi-user mode (value {\tt enabled}) or in single-user model (value {\tt disabled}).
By default, the multi-user mode is enabled, allowing multiple users with different user
names to connect to the CBserver. In single-user mode, only clients started
by the same user (identified by her name) can connect to the CBserver.
The single-user mode is recommended if you want to block other users from logging
into your CBserver. Since the test is done only on the user name, a malicious
attacker could use your user name for an account on another computer and then
successfully log into your CBserver. Use Internet firewalls to protect
against such attacks. If you specify an administrator user (option -a), then
this use can always connect to the CBserver.

\item[-v vmode] controls whether view maintenance rules are generated (vmode={\tt on})
or not (vmode={\tt off}). View maintenance rules are used to keep a ConceptBase.cc
view up-to-date upon changes to the object base. Default value for vmode is {\tt off}.

\item[-mc maxcost] this parameter defines the maximum cost level for a predicate in a 
binding path that is used to compile a meta formula (see section \ref{sec:CCmeta}). 
The evaulation of a binding path yields fillers for the meta variables. Set maxcost to
10 if such a predicate should have about one free variable. Set it to 100 if
if may have two free variables. Default is 100. The higher the number, the more
candidate paths are generated, increasing the likelihood that a binding path is 
found. On the downside, a high value increases the compile time of meta formulas.

\item[-pl pathlen] sets a maximum length for binding path candidates. In principle,
the number of candidates can explode with the path length. Like the previous parameter,
the path length influences the ability of ConceptBase.cc to compile meta formulas. The
default value is 5. If you set the value to 0, then no meta formula can be compiled.

\item[-im imax] sets the maximum number of iterations used to re-order attribution 
predicates with one free variable. Evaluating such predicates first can lead
to faster elimination of free variables and thus lead to better query and ECA
performance. The default value is 3.

\item[-eca emode] controls the ECA sub-system. Possible values are {\tt unsafe}
(ECArules are evaluated without safeguarding
recursive deductive rules), {\tt off}
(ECArules are not evaluated, even if some are defined), and {\tt safe}
(ECA rules are evaluated with safeguarding recursive rules; this is the default).
Use the mode {\tt unsafe} if none of your ECArules calls recursive predicates
on the newest database state. This may lead to a limited speed-up.

\item[-eo eomode] controls the optimization of conditions of ECArules.
Possible values are {\tt on} (default) and {\tt off}. The optimization is done by
a re-ordering of predicates in the condition. Hence, you only want to turn of
optimization to gain full control over the order of evalution in ECA conditions.

\item[-load dir] specifies the directory from which the CBserver will load
module sources at start-up time. The module sources must have file names
starting with {\tt System} and file type {\tt sml}. Typically, they are
generated via the {\tt -save} flag in the preceding session of the CBserver.
The default is {\tt none}, i.e. no module sources are loaded at CBserver
start-up.

\item[-save dir] specifies the directory into which to save certain textual
excerpts of the database, in particular module listings. The parameter has the
default value {\tt none}, which disables the saving function. Currently,
the CBserver only saves module listings. Each module is stored in one
file with file type {\tt sml}. The directory {\tt dir} must exist.
The module listing is performed when the CBserver is shut down (complete module
tree is listed), or when a client tool logs out (home directory tree of
the client tool is listed). The module listings uses the {\tt set module}
directive to enable the import of the file to the right module location.
See also section \ref{sec:module_sources} for details.


\item[-views dir] specifies the directory into which the results of certain
queries are materialized. See section \ref{sec:module_views} for details.

\item[-ms sep] specifies the module separator to be used for saving module listings
and views. If the separator is set to '-'  (default), then all module sources and
views are stored at the top level directory. If the module separator is set to '/',
then the files are stored in a deep sub-directory structure that mirrors the module
structure.

\item[-mg mgmode] specifies whether module listings are generated with separators
"{\tt \{---\}}" for each transaction occurring in the module (option {\tt split}) or
without such separators
(option {\tt whole}). The default is {\tt split}. The split option better supports
cases where metaformulas are defined and used in the same module. The third option
is {\tt minsplit}. It minimizes the number of separators to those that are essentially needed.
Subsequent loading of such module sources is then faster.

\item[-rl rlmode] controls the way how the CBserver creates labels for generated formulas.
The default value is {\tt on}, instructing the CBserver to find a readable 
label for the generated formula. It typically consists of the labels of
the participating metaclass attributes occuring in the metaformula. 
If set to {\tt off}, the CBserver will just take a system-generated label that contains
a unique identifier. This is slightly less readable (if you want to inspect
the generated formulas) but safe against certain possibilities of assigning
the same label twice.

\item[-ia tmax] sets the maximum number of hours during which a client should
interact with the CBserver to be regarded as active. Negative values
are interpreted as 'infinity'. This parameter is only used when a CBserver
uses the -sm slave and -r options, or the -g public option.
The default value for tmax is 2.0 hours.

\item[-sm servermode] sets the server mode. Possible values are
{\tt master} (default) and {\tt slave}. In slave mode, the last client that
leaves the CBserver will also shutdown the CBserver, provided that the CBserver
and the client were started by the same user. This option is useful when
a CBserver is only needed while still clients are registered. A master
CBserver must always be stopped explicitly.

\item[-st stratmode] enables or disables the rule stratification test. Possible values are
{\tt on} (default) and {\tt off}. If enables, then the query evaluator shall
dynamically test whether stratification violations occur. They shall then
be reported as an error. Disable the test, if you are sure that the answers
are correct even though a stratification violation occurs.

\item[-g cmd] provides a special command to the CBserver. There are currently
three such commands. The command {\tt nolpi} instructs the CBserver to ignore
any plug-in file (see section \ref{cap:lpi}). The command {\tt public} configures
the CBserver as a public CBserver (see \ref{sec:pubcbserver}).
The command
{\tt exit} instructs the CBserver to exit immediately after start-up. This can
be useful to combination with the option {\tt -views}, {\tt -db} and {\tt -save}
to materialize some excerpts from a stored database.

\item[-a user] designates {\tt user} as 'administrator' of this CBserver. 
For the time being this just gives the right to shutdown the CBserver. By default,
the user who started the CBserver is its administrator. This user shall
also keep the right to shutdown the server, even when another user is the
designated administrator. If you specify the user name with host, e.g.\ 
{\tt billy@myhost}, then only the user {\tt billy} on host {\tt myhost} is
recognized as additional administrator. 


\end{description}


\vspace{0.5cm}

If a CBserver is started without any parameter, then the update mode shall be set
to {\tt nonpersistent}, the trace mode to {\tt no}, multi-user mode is disabled,
and the server mode to {\tt slave}. The other parameters are set to their defaults.
Such a CBserver is useful as companion of tools that need it only while they
are running. 

\begin{verbatim}
   cbserver
\end{verbatim}

A ConceptBase client running on the same computer will then connect to this
CBserver, when it uses 'localhost' as host and 4001 as port number. Since
the CBserver runs un slave mode, it will shut down when its client disconnects.

\section{ConceptBase under Windows 10}
\label{sec:win10}
The CBserver is only compiled for Linux architectures. This means that user of
other platforms must rely on a Linux system to utilize ConceptBase. The traditional
way is to start the CBserver on such a Linux system and then connect to it, possibly
using a so-called public CBserver (see section \ref{sec:pubcbserver}).
Since April 2017, this detour is no longer required for users of Windows 10 (64bit,
Creators Update). This version of Windows is capable to let Linux programs
run under the 'bash' utility, which is basically a whole Linux system under Windows
that realizes calls to the Linux API by hooks to the Windows API. Hence, it is not
a virtual machine, it lets you run the Linux (64bit) variant of the CBserver
natively on Windows 10.

To enable the Linux capability on Windows 10, follow the instructions at
\url{http://conceptbase.sourceforge.net/CB-WinLinux.html}. Note that you must have
installed Java (64bit) on your Windows machine, not Java (32bit) to take full advantage
of this feature. You can check whether your Java is 64bit by calling the following command
in a Windows command window.

\begin{verbatim}
   java -d64 -version
\end{verbatim}


Users of older Windows versions and of other operating systems can continue to use the
public CBserver (see section \ref{sec:pubcbserver}) to take advantage of ConceptBase.



\section{Database format}

A ConceptBase database is a directory that contains at least the following files:

\begin{itemize}
\item OB.symbol: a binary file that associates object names (like 'MyClass') with
object identifiers. 
\item OB.telos: binary file storing all propositions
\item OB.rule: text file containing the generated Prolog code for rules, constraints,
and queries
\item OB.ruleinfo: text file containing argument information about queries and some
formation for the cost-based formula optimizer
\item OB.ecarule: text file containing the generated Prolog code for active rules
\end{itemize}

The database files may only be updated via the ConceptBase server. Their
initial state is bootstrapped from textual Telos frames that define the
pre-defined objects of O-Telos.  Since the pre-defined objects can change
from version to version, we cannot guarantee binary compatibility of 
ConceptBase databases. You can easily export the textual definitions from
a databases via the -save option. Those definitions can then be 
imported to the new ConceptBase version.
The database directory may contain further text files with filetype 'lpi'. These are
Prolog plugins loaded at startup time, see also section \ref{cap:lpi}.


\section{Modifying the system database}

A new database is created from the database lib/SystemDB in your ConceptBase
installation directory. The System database contains exactly the objects
of the root module {\tt System}. They include for example the definitions
of the objects {\tt Proposition}, {\tt Individual},  {\tt Attribute},
 {\tt InstanceOf}, and  {\tt IsA}. Further the objects {\tt Class},
 {\tt QueryClass},  {\tt Function},  {\tt ECArule},  {\tt Module}
and many more are defined that are needed to formulate queries and
to use the capabilities of the system. 

Whenever a new database is created, the files from
this System database are copied into the new database directory. This allows experienced
ConceptBase users to adapt the System database to their needs. Just start a 
ConceptBase server with

\begin{verbatim}
cbserver -d $CB_HOME/lib/SystemDB -s 0
\end{verbatim}

Replace {\tt \$CB\_HOME} by the path to your ConceptBase installation directory.
Then start a CBIva user interface, connect to the CBserver and switch to the
module {\tt System}. Assume you want to predefine the class {\tt Container}
and declare a {\tt Model} as subclass of {\tt Container}:

\begin{verbatim}
Container with
  attribute
     contains: Proposition
end
Model isA Container end
\end{verbatim}

You can also add rules and constraints about containers, e.g. that containers
may not contain themselves.
A more significant extension would be to add active rules to the system database.
For example, the active rules
in CB-Forum at \url{http://merkur.informatik.rwth-aachen.de/pub/bscw.cgi/3260276}
changes the semantics of the UNTELL operation.
Such definitions are subsequently included in any new database that you create.
Be careful with deleting system objects. The code of the CBserver relies on
the existence of certain system objects. 


\section{Tracing and restarting}

\vspace{0.5cm}
The trace of the CBserver can be saved by redirecting its
output, e.g.

\begin{verbatim}
cbserver -r 10 -port 4444 -t high -d MYDB >> mylogfile.log
\end{verbatim}

The CBserver can also be started directly from the ConceptBase.cc
User Interface (see section \ref{cha:workbench}) and most parameters can
be specified interactively. The command line version is recommended when
one CBserver serves multiple users or when user interface and server shall run
on different machines.
The parameter {\tt -r 10} instructs the CBserver to restart after 10 seconds if 
a crash has occurred.




A special error message during the startup of the CBserver
is the following:
\begin{verbatim}
### FATAL ERROR:
Application is locked by hostname, PID 1234 
### CBserver aborted
\end{verbatim}
This messsage is printed if there is still a file with the name OB.lock
in the database directory (option -d). The OB.lock file should avoid that
two servers are using the same database directory. The file may be left over of a previous
CBserver if the server was not stopped correctly (e.g. aborted by Ctrl-C or it crashed).
If you get this error message, make sure that there is no other server running
that uses this directory and then delete the file OB.lock. Then, the 
CBserver should start correctly.


\section{Public CBservers}
\label{sec:pubcbserver}

A public CBserver is a ConceptBase server that is accessible from 
the whole network. As such this is a property that any CBserver has, except when
the multi-user capability is disabled, or when your firewall prevents
external access.

If you work with an existing database, you may want to specify some access rights rules like suggested in section \ref{sec:modaccess}.
We neglect them in this simple example.
Now, start the public CBserver with suitable parameters under Linux/Unix:

{\small
\begin{verbatim}
cbserver -r 2 -a jonny -g public -ia 0.5 -u nonpersistent -d MDB &> log.txt &
\end{verbatim}
}

The option {\tt -g public} enables the slave mode implicitely and tells {\tt oHome}
as an instance of {\tt Auto\-Home\-Module}. The combination of the slave mode and the
option {\tt -r} instructs CBserver to stop and
restart when the last client exits.
The update mode is nonpersistent. As a consequence, the restarted CBserver
will use the unchanged state of {\tt MDB}. This is useful, if you want to provide
the service of CBserver to a larger group of (anonymous) users. They can log in
with a client, operate on the database in nonpersistent mode and eventually leave
the CBserver. When the last active client\footnote{An active client is a client 
whose last interaction with the CBserver was less than a certain
number of hours ago, specified with the CBserver parameter {\tt -ia}.}
leaves the CBserver, then CBserver will freshly start up
after 2 seconds. Due to the option {\tt -u nonpersistent}
user-defined objects are only stored at the public CBserver
while there are still active clients enrolled to the public CBserver. 

The parameter {\tt -a} sets the administrator user of the public CBserver. This user
is allowed to shutdown the CBserver from a client.
The auto home feature will assign different users to individual workspaces.
Unless you introduce access rights (and enable them via the CBserver option {\tt -s 2})
the users can also manipulate the modules of other users. However, the CBserver is
restarted whenever the last client logs off. Hence, the definitions of different users
are not permanently stored on the CBserver.
The parameter {\tt -ia} used here instructs the CBserver to regard a client as active of it
had its last interaction within 0.5 hours. Clients that were inactive for a longer
time will not prevent the CBserver from restarting when another (active) client 
logs off the CBserver.

Consider the following alternative command to start a public CBserver:

{\small
\begin{verbatim}
cbserver -t no -r 2 -a jonny -g public -ia -1 -d ALLDB &> /dev/null &
\end{verbatim}
}

Here, the updates are stored persistently in the database ALLDB. Changes to modules are not lost
when the last active client leaves the CBserver. There is no log file created as well.
The option '-ia -1' used here instructs the CBserver to regard any client as active,
regardless of how long ago its last interaction occurred.


Configure the CBIva interface to use the public CBserver via the variable {\tt PublicCBserver}.
It can be set by the menu item "Options/Edit Options Manually" of CBIva, see section \ref{sec:config}.
A value different from {\tt none} enables the use of the public CBserver by CBGraph.
You can optionally append a port number like "cbserver.acme.com/4002". The default value
for the port number is 4001. Installations of ConceptBase on platforms, for which no binaries of the CBserver
exist, may use a default public CBserver. This is the case OS-X and older versions of Windows%
\footnote{Windows 10 (64bit) Creators Update has the ability to activate a Linux sub-system
which can then be used to start CBserver transparently.
If you had previously used a public CBserver under Windows 10 and now want to utilize a local CBserver,
then reset the variable {\tt PublicCBserver} to {\tt none}.}.

Don't forget to save the options after changing the variable.


The best way to interact with a public CBserver is to use graph files, see section
\ref{sec:cbgraphcmd}. Assume that a public CBserver is running on host {\tt cbserver.acme.com}.
and that {\tt cbserver.acme.com} was configured as the public CBserver to be used.
Then, calling

\begin{verbatim}
cbgraph graph1.gel 
\end{verbatim}

will connect to the public CBserver instead of 'localhost'. The port number for the connection is taken
from the graph file.
Do not forget to save the graph file before exiting CBGraph if the public CBserver was started
in non-persistent mode.
If you subsequently open the graph file, it will attempt to connect to the same host.
You can force it to attempt the connection to localhost instead by

\begin{verbatim}
cbgraph -host localhost graph1.gel 
\end{verbatim}

If you are using ConceptBase under Linux, then the CBserver is running by default on your local computer
("localhost"). The same is true for Windows 10 with an enabled Linux subsystem. Users on Mac
computers or older Windows versions are by default using a public CBserver at the university where
ConceptBase is developed. This public server is meant for testing the system. {\bf DO NOT USE it for
managing confidential information!} If you plan to use ConceptBase in a serious way, we
recommend that you set up a protected Linux server on your own network that runs the CBserver.
Users can set up the variable "PublicCBserver" in CBIva to the address of the protected Linux server
to automatically connect to that server, see section \ref{sec:config}. You can also connect to that
server manually via the "Connect" function of the ConceptBase clients CBIva, CBGraph, and CBShell.




\section{The tabling subsystem}
\label{sec:cbcache}

Since version V5.2.4 ConceptBase.cc features a new query evaluation method,
which uses a so-called tabling cache to store intermediate results of predicates
that are called during the top-down (SLDNF) query evaluation. Assume, for example, that
an employee 'bill' has two projects 'p1' and 'p2'. Then, the
result of a predicate '(bill hasProject x)' with variable x would be
the set \{(bill hasProject p1),(bill hasProject p2)\} consisting of facts. 
We call this fact set also the {\em extension} of the predicate.

After a completed predicate evaluation, the tabling cache of the predicate holds
its extension. Tabling speeds up query evaluation and
prevents infinite loops when ConceptBase.cc evaluates
recursive queries and deductive rules. Essentially, the
tabled evaluation allows to compute dynamically stratified
semantics of the Datalog database underlying ConceptBase.cc.
Plenty of examples for recursive rules and queries are provided
in the online ConceptBase Forum.

The CBserver provides three tabling cache modes to control the
behavior during query evaluation:

\begin{description}

\item[-c off] In this mode, the cache is completely disabled. Use this mode
when your models do not include recursive rules.
The mode is only provided for backward compatibility and has no advantages.

\item[-c keep] The cache is only emptied when necessary, in particular when the cache
has been invalidated by an update to the database, or when the maximum number
of facts in the cache is exceeded. The maximum number is currently set
to the default 20000. Exceeding the maximum is not an error. It only indicates that the cache
is marked for being emptied. If necessary, the cache emptying takes place before
a transaction. The keep mode is on average consuming more main memory than the
transient mode but speeds up response time enormously in case of re-use of query
results. The 'keep' mode is the default mode for tabling. You can change the maximum
cache size by the CBserver command line option "-cs".

\item[-c transient] The tabling cache is emptied before
each transaction (ask, tell, untell, retell). A subsequent query is always
evaluated starting with an empty cache. This mode is somewhat 'safer' than the
'keep' mode since it starts each query with an empty cache state. 
While the answer to a query is in principal independent from the cache mode,
the cache mode has a certain influence on the persistence of objects created
within a transaction. Specifically, results of arithmetic expressions 
computed during one transaction shall be removed after the transaction
when the cache mode is 'transient'. In cache mode 'keep', these objects
continue to exist and are visible to future transactions. 

\end{description}

ConceptBase will only call tabled evaluation for deductive predicates.
Other predicates are evaluated by the regular SLDNF engine of the
underlying Prolog engine.
By default, the tabling cache mode is activated in mode {\tt keep}.
Some statistics on cache usage are written to the terminal window of
the ConceptBase server when the tracemode has been set to {\tt veryhigh}. 

Acknowledgements: The techniques for the
tabled query evaluator of ConceptBase.cc are inspired by the 'tabled evaluation' \cite{ssw94,chenwarren96}.
We do however not delay the evaluation of negated predicates but rather re-order
them at compile time to guarantee that all variables are bound at call time.
Tabled evaluation is also implemented in XSB [\url{http://xsb.sourceforge.net/}]
and DES [\url{http://www.fdi.ucm.es/profesor/fernan/des/}].


\section{Database persistency}
\label{sec:persistency}

The default update mode is 'persistent'. In persistent mode, all changes
to the database are written to the file system at the directory specified in
the parameter '-d'. Persistent mode is suitable when a CBserver
runs for a longer period of time and is directly updated by
application programs. 
If ConceptBase.cc is used for testing and modeling
purposes, the update mode 'nonpersistent' is an interesting alternative.
We discuss two scenarios for the nonpersistent mode and one for the
persistent mode.
\newline

{\bf Scenario 1: Single-user modeling}. When a user needs to model a certain
application domain with classes and meta classes, he usually works
with external Telos files (aka source models, file extension '.sml').
These files can include comments like usual with program source code.
The recommended mode here is '-u nonpersistent' without specifying a
database. The user can load the source models
into such a non-persistent server and make corrections to the source
files in case of errors or design changes. Here, ConceptBase is
mostly used to check and analyze the models.
Recommended options:
\begin{verbatim}
cbserver -u nonpersistent -mu disabled
\end{verbatim}

{\bf Scenario 2: Lab assignments}.
Assume that a teacher wants students to exercise a certain modelling task
using ConceptBase.cc. Then, he would prepare some Telos files with necessary
definitions (e.g. some meta classes) and load them into a persistent
ConceptBase server. After that, he can restart the ConceptBase server
in non-persistent mode on the same database created before. Student
can then work on their extensions while the state of the database
can easily be set back to the original state defined by the teacher.
The module system of ConceptBase.cc can be used to support multiple
students to work on the same server without interfering with each other,
see section \ref{sec:module}.
Recommended options:
\begin{verbatim}
cbserver -d MYDB -s 2 -mu enabled -u nonpersistent
\end{verbatim}

The second scenario might also be useful in modeling. If there are some
parts that are regarded as stable, the modeller can decide to make
them persistent and only add/modify those Telos models that are still
subject to change. In particular for large Telos models, this strategy
saves time.
Note that the updates by the users are lost when the non-persistent
CBserver is stopped. This might be useful, if you want to re-use the same
initial database MYDB several times, e.g. for different user groups.
\newline

{\bf Scenario 3: Project work}.
Here the students work for several days on a given task. Changes shall 
not be lost. The use of the {\tt -db} option will not only store the database
in binary form but also store the source code of all modules as text
files in directory MYDB.
Recommended options:
\begin{verbatim}
cbserver -db MYDB -s 2 -mu enabled 
\end{verbatim}


If ConceptBase.cc is used in a multi-user setting, then one can combine
the update mode with the module feature (see section \ref{sec:module}).
In this scenario, multiple users access the same CBserver.
A common super module (e.g. the module {\tt oHome}) carries the common
objects of the users. Each user can be assigned to her own hown module
(a sub module of the common super module) and create and update 
objects in this workspace without interfering with other users.
If several groups of users shall share their definitions, then they would
be assigned to the same home module. The home module may have sub modules
for testing and releasing definitions. By employing access rights to
modules, one can also design which user has which read/write permissions.
The builtin query {\tt listModule} allows to save the contents of a
module to a Telos source file (see section \ref{sec:module_content}).






\section{The UNTELL operation}
\label{sec:untell}

ConceptBase.cc realizes the concept of a historical database. The TELL operation
submits O-Telos frames to the CBserver. The CBserver extracts the 'novelty'
of the submitted frames and translates it into a set of P-facts to be stored
in the object store. Any P-fact has a so-called {\em belief time\/} associated to
it (see section \ref{sec:representation}). The belief time is an interval $(t_1,t_2)$ whose left boundary $t_1$ is the time point when
the P-fact was inserted to the object store, i.e.\ the time when the transaction
was executed that led to the insertion of the P-fact. The right boundary $t_2$ specifies
the time point after which the CBserver assumes the P-fact to be not true anymore.
It is initialized with 'Now' when the P-fact is created. This symbolic value
is interpreted as the current time. You may also interpret such a time interval to be right
open.

The UNTELL operation terminates the belief time of P-facts specified in an O-Telos
frame. The value
'Now' is replaced by the time%
\footnote{ConceptBase takes the system time using timezone Coordinated Universal Time (UTC) of the computer on
which it is running and rounds it to milliseconds. The time is captured when the transaction is initiated,
i.e. all P-facts told or untold in the transaction will use that transaction time.
}
when the UNTELL operation is executed. 

From the user's  perspective, a TELL operation is about creating some objects and
the UNTELL operation is about deleting them%
\footnote{ 'Deleted' objects can however be recovered
by setting the so-called roll-back time before an ASK transaction is issued.
Only ASK transaction are allowed on historical database state. It makes little
sense to update a historical state. That would be 'falsifying' the history of stored P-facts.}%
. Many users expect the UNTELL operation to be symmetric to the TELL operation,
i.e.\ untelling a frame that has been told before should remove the frame completely. 
This is however not the case for the following reasons:

\begin{itemize}
\item An O-Telos frame being argument of a TELL or UNTELL operation is not necessarily all
the information about an object but just some. 
\item Other objects might refer to an object told previously. An UNTELL operation would then
be rejected to preserve referential integrity.
\item ConceptBase.cc adds instantiation to the builtin classes {\tt Individual},
{\tt Attribute}, {\tt Instance\-Of}, and {\tt IsA} depending on the type of the object.
\end{itemize}

The last reason is most significant in preventing symmetry. As an example consider the O-Telos 
frame (referred to as frame 1).

\begin{verbatim}
  bill in Employee with
    name
     bname: "William"
  end
\end{verbatim}

ConceptBase.cc will recognize that {\tt bill} is an individual and that {\tt bill!bname} is an
attribute. This information is attached internally to the P-facts, more precisely, it is derivable
from the structure of the corresponding P-facts.
Assume that you started a CBserver with untell model {\tt verbatim} (see below).
When you ASK the CBserver for the frame of {\tt bill}
after the TELL operation on frame 1, you will get frame 2:

\begin{verbatim}
  Individual bill in Employee with
    attribute,name
     bname: "William"
  end
\end{verbatim}

Hence, the instantiation to the builtin classes {\tt Individual} and 
{\tt Attribute} is added to the frame. If we submit the original
frame 1 to an UNTELL operation, ConceptBase.cc assumes by default that only two facts should be untold:

\begin{enumerate}
\item The instantiation of {\tt bill} to {\tt Employee}.
\item The instantiation of the attribute {\tt bill!bname} to its attribute category {\tt Employee!name}.
\end{enumerate}

As a consequence, the object {\tt bill} and its attribute continue to exist after the UNTELL on frame 1.
It would look like (frame 3):

\begin{verbatim}
  Individual bill with
    attribute
     bname: "William"
  end
\end{verbatim}

Only by untelling this frame 3 as a second operation or by untelling the completed frame 2, the
object {\tt bill} and its attribute {\tt bill!bname} shall be made historical.

This assymmetry of TELL and UNTELL is regarded by some ConceptBase.cc users as unnatural behavior.
They expect that an UNTELL is also supposed to affect the objects themselves, not just their
instantiation to classes. To support those users, we provide a so-called untell mode (see parameter
-U in the list of CBserver command line parameters). The untell mode {\tt verbatim} will cause
ConceptBase.cc to behave as explained above. 
The mode {\tt cleanup} (default) will take care to remove the objects themselves
provided that they have no other instantiations to classes except instantiation to the four
builtin classes {\tt Individual}, {\tt Attribute}, {\tt InstanceOf}, and {\tt IsA}. Furthermore,
no other object may be linked to the object subject to be untold. 

\subsection{Cascading UNTELL}
\label{sec:cascuntell}

You can use ECArules to realize a cascading UNTELL, i.e.\  if an object is deleted then
all its links to other objects are deleted as well, leading to potential follow-up deletions.
The required ECArules are provided in 
\url{http://merkur.informatik.rwth-aachen.de/pub/bscw.cgi/3260276}.
You only need to add the ECArules to your database to enact the cascading UNTELL.
You should be careful with using the cascading UNTELL. For example, untelling
a system class will also untell all instantiations to it. Use the CBserver
option {\tt -s 1} to prevent such undesired deletions.

The model CascUntell.sml contains only ECArules to complete UNTELL operations. This makes
it less cumbersome to remove objects or attributes of objects. The model CascUntell-Plus.sml
adds ECArules that check also the presence of TELL operations to attributes (including
links to other objects). If you add these rules to your database, then you can overwrite
attributes just by telling the new frame. For example, assume that "mary" was initially defined
as

\begin{verbatim}
  mary in Person with 
     name marysname : "Mary"
     age marysage : 32
   end 
\end{verbatim}

Some time later, you want to revise the object as follows:

\begin{verbatim}
  mary in Person with 
     name marysname : "Mary"
     age marysage : 33
   end 
\end{verbatim}



The ECArules in CascUntell-Plus.sml will detect that an attribute "mary!marysage" has already
an old value "32". Consequently, the old attribute and its categories such as "age" are removed
and the new attribute is inserted.
Without the ECArules, you would have to UNTELL the old attribute manually and then tell the new
attribute. Note that the service of the ECArules is only active if the ECArules are stored in the
database.




\section{Memory consumption and performance}
\label{sec:restrictions}

ConceptBase.cc stores objects in a dedicated object store maintained in
main memory. A P-fact $P(o,x,n,y)$ consumes about 800 bytes of main memory.
That means that one can store roughly 1 million P-facts in 1 GB of main memory.
A typical Telos frame is stored with roughly 10 P-facts. Hence, 1 GB of main
memory allows you to store around 100.000 Telos frames.
On 32 bit CPUs, this results in a maximum of roughly 400.000 frames that
can fit into 4 GB of addressable main memory. This restriction virtually
vanishes with 64 bit CPUs.

A single TELL/UNTELL operation submitted to the CBserver should
not contain more than about 2000 frames (at about 5 attributes per frame).
Otherwise, the compiler can run out of stack memory.

The raw performance of the object store, i.e.\ the time needed to reconstruct a 
frame for a given object identifier, is virtually independent from
the number of P-facts that it stores. However, if you have defined many rules
or integrity constraints, the performance may well degrade significantly
with the number of stored P-facts. The same holds for queries. 
We tested the response times for standard queries such as computing the
transitive closure in relation to varying database sizes. Results indicate that
ConceptBase apparently approximates in many cases the theoretic optimum. 

The performance of the active rule evaluator (section \ref{sec:eca}) is currently 
rather limited. We measured around 100 rule firings per second. This can be a 
performance bottleneck when many active rules are being processed.



\section{The Java API to the CBserver}
\label{sec:javaapi}

The application programming interface (API)
to the ConceptBase server (CBserver) is realized by the Java class 
{\tt LocalCBclient}.
Most (Java) application programmers presumably only need the simple String-based
part of {\tt LocalCBclient} to interact with a CBserver. {\tt LocalCBclient} uses
socket-based data streams to realize the communication.
We define here the methods of this String-based API. If the argument starts with an "s",
then the data type is {\tt String}. If it starts with an "i", then the data type
is {\tt int}.


\begin{verbatim}
LocalCBclient cbClient = new LocalCBclient();
\end{verbatim}

\begin{itemize}
\item This constructor provides the API object {\tt cbClient} to be used for the
subsequent method calls.
\end{itemize}

\begin{verbatim}
sAnswer = cbClient.cbserver();
\end{verbatim}

\begin{itemize}
\item This method attempts to start a single-user CBserver in "slave mode" 
on localhost with port number 4001.
The method only works on platforms for which the CBserver was compiled. This is currently
limited to Linux and Windows 10 (see also section \ref{sec:win10}). In other cases,
you need to start a CBserver on a host system that supports the CBserver and then use the 'connect'
method. If the CBserver was successfully started, the method returns "yes", otherwise "no".
If the cbClient is already connected to a CBserver, then the method also returns "yes".
The CBserver shall automatically shut-down when the last client disconnects from it.
\end{itemize}

\begin{verbatim}
sAnswer = cbClient.connect(sHost, iPort, sTool, sUser);
\end{verbatim}

\begin{itemize}
\item This method connects your Java program to the CBserver specified by {\tt sHost} (the domain
name of the computer on which your CBserver runs on) and {\tt iPort} (the port number is 
an integer; usually 4001). The string sTool shall be a self-selected name of your
tool, e.g. {\tt "ModelerXY"} and sUser is a string containing a user name (typically your
user name). If you use {\tt null} as username, then the {\tt cbClient} will use the
login name of the computer user that started the Java program.
The return value is {\tt "yes"} if successful and {\tt "no"} else. Use the boolean-valued
function {\tt cbClient.isConnected()} to check whether {\tt cbClient} is currently
connected to a CBserver. If the socket connection breaks down, then the {\tt cbClient} will
set the connection status to unconnected.
\end{itemize}

\begin{verbatim}
sAnswer = cbClient.connect();
\end{verbatim}

\begin{itemize}
\item Connects to the public CBserver  (see section \ref{sec:pubcbserver})
if configured for your installation.
Otherwise, attempts to start a single-user "slave" local CBserver on port 4001 if not started.
Connects to the local CBserver with the default tool name "LocalCBClient" and {\tt null} as username.
It behaves similar to the "connect" command of CBShell.
\end{itemize}


\begin{verbatim}
sAnswer = cbClient.disconnect();
\end{verbatim}
\begin{itemize}
\item This method disconnects connects your Java program from the CBserver.
The return value is {\tt "yes"} if successful and {\tt "no"} else. If the CBserver was started in
"slave mode" and the cbClient is the last remaining client of the CBserver, it will shut down.
\end{itemize}

\begin{verbatim}
sAnswer = cbClient.pwd();
\end{verbatim}
\begin{itemize}
\item This method outputs your current working module path, e.g. {\tt "System-oHome"}.
\end{itemize}

\begin{verbatim}
sAnswer = cbClient.mkdir(sModule);
\end{verbatim}
\begin{itemize}
\item This method creates a new submodude (e.g. {\tt "MyMod"}) in the current module.
\end{itemize}

\begin{verbatim}
sAnswer = cbClient.cd(sModule);
\end{verbatim}
\begin{itemize}
\item This method changes the current module to {\tt sModule}, e.g. {\tt "MyMod"}.
\end{itemize}

\begin{verbatim}
sAnswer = cbClient.tells(sFrames);
\end{verbatim}
\begin{itemize}
\item This method tells (=defines) the objects given by the string sFrames to the
current module of the CBserver. 
The return value is {\tt "yes"} if successful and a string containing user-readable error messages else.
\end{itemize}

\begin{verbatim}
sAnswer = cbClient.untells(sFrames);
\end{verbatim}
\begin{itemize}
\item This method untells (=removes) the objects given by the string {\tt sFrames} from the
current module of the CBserver. 
The return value is {\tt "yes"} if successful and a string containing user-readable error messages else.
\end{itemize}

\begin{verbatim}
sAnswer = cbClient.asks(sQuery,sFormat);
\end{verbatim}
\begin{itemize}
\item This method asks the query call {\tt sQuery} (given by a String) to the CBserver 
The return value is a string containing the answer to the query.
The return value is {\tt "no"} in case of errors, e.g.\ when the query is not defined.
The query call can have arguments, e.g. {\tt "get\_object[Class/objname]"}. 
The CBserver shall use the answer format {\tt sFormat} for the result. Thus,
if you want to define your own answer format, then use the facilities
described in section \ref{sec:answerformat}. The query shall be answered
in the context of the current module and the current time ({\tt "Now"}).
The answer {\tt "nil"} stands for an empty answer.
\end{itemize}

\begin{verbatim}
sAnswer = cbClient.asks(sQuery);
\end{verbatim}
\begin{itemize}
\item Same as {\tt cbClient.asks(sQuery,"default")}, i.e. the CBserver determines the
applicable answer format (in most cases: {\tt "FRAME"}).
\end{itemize}

\begin{verbatim}
sAnswer = cbClient.clientid();
\end{verbatim}
\begin{itemize}
\item Return the identifier by which this client is registered to a CBserver.
\end{itemize}

\begin{verbatim}
sAnswer = cbClient.clearall();
\end{verbatim}
\begin{itemize}
\item Attempt to delete all objects from the current module. Returns "yes" if
successfull, otherwise "no". Calls the 'purgeModule' query described in section
\ref{sec:module_purge}. Use this method with great care since it wipes out
the whole module content.
\end{itemize}



\vspace{0.5cm}
Below is the listing of the Java program {\tt TinyClient.java} that uses the 
String-valued API:

{\small
\begin{verbatim}
import i5.cb.api.*;
public class TinyClient {
  private static LocalCBclient cbClient = null;
  public static void main(String argv[])  {
    String sAnswer;
    cbClient=new LocalCBclient();
    sAnswer = cbClient.connect("cbserver.acme.org",4001,"TinyClient",null);
    sAnswer = cbClient.tells("Employee in Class end");
    sAnswer = cbClient.asks("get_object[Employee/objname]");
    System.out.println(sAnswer);
    sAnswer = cbClient.disconnect();
  }
}
\end{verbatim}
}

You need to compile the Java program with the {\tt cb.jar} library. This is available 
from the ConceptBase installation directory (referred here as CB\_HOME). To compile the Java
program call

\begin{verbatim}
  javac -classpath $CB_HOME/lib/classes/cb.jar TinyClient.java
\end{verbatim}

Before running the client, make sure that the CBserver runs on the specified Linux computer
(here "cbserver.acme.org") under the specified port number (here 4001), and that
this port number is accessible from your client computer.
You may want to configure that CBserver as a public CBserver (see section \ref{sec:pubcbserver})
to have a save way to connect your Java program to it.

The Java program can then be started under Linux as follows:

\begin{verbatim}
  java -classpath $CB_HOME/lib/classes/cb.jar:. TinyClient
\end{verbatim}

Under Windows, you should use:

\begin{verbatim}
  java -classpath c:\conceptbase\lib\classes\cb.jar;. TinyClient
\end{verbatim}

The Java source code for {\tt TinyClient.java}, {\tt TinyClient2.java} (uses the cbserver method),
{\tt TinyClient3.java} (uses the connect method),
 and a more elaborate example 
{\tt SimpleClient.java} is included in the directory {\tt examples/Clients/JavaClient}
of your ConceptBase installation directory.


\chapter{The CBShell Utility}
\label{cha:cbshell}

The ConceptBase.cc Shell (CBShell) is a command line client for ConceptBase.cc. It allows
to interact with a CBserver via a text-based command shell.
Moreover, it can process commands from a script file without further user
interaction. The utility can be employed to automate certain activities such
as batch-loading a large number of Telos models into a CBserver,
or to extract certain answers from a CBserver.

\section{Syntax}

There are two ways to use the CBShell. The first one processes the commands from a script
file (batch mode). The second one prompts the user for commands (interactive mode).

\begin{verbatim}
   cbshell [options] scriptFile [params]
   cbshell [options]
\end{verbatim}


\section{Options} 


\begin{description}

\item[-l] This options instructs CBShell to write errors and some statistic information
to the files error.log and stat.log.

\item[-f scriptFile] Execute the commands specified in scriptFile rather that requesting commands
from the command line interface. If the -f option is used and a scriptFile
is specified, the commands of the file will be executed without user interaction,
and CBShell will exit at the end. The prefix -f can also be omitted, i.e.\
"cbshell scriptFile" is equivalent to "cbshell -f scriptFile".

\item[-t] This option can only be used in combination with the -f option. It shall instruct CBShell
to confirm each command in the script file before it is executed.

\item[-i] This option modifies the 'cbserver' command by invoking a CBserver compiled directly from its sources.
For developers only.

\item[-a] This options instructs CBShell not to directly show the answer to
   each command in interactive mode;  instead you have to manually call showAnswer

\item[-v] This options enables the verbose mode. In this mode, the command and the answer
   are always displayed on standard output

\item[-p] Disables the display of the command prompt in interactive mode. This may be
useful when CBShell is used in a Unix pipe where the preceding program generates the
commands and feeds them into CBShell.

\item[-q] Instructs CBShell to convert single quotes in positional parameters into escaped double quotes.
This option is useful when a parameter contains special characters and still shall be regarded as a valid
object label by ConceptBase. Useful when calling CBShell scripts within regular scripts (e.g. bash)
that pass parameters with special characters to the CBShell scripts.
See also section \ref{sec:shellintegration}.

\item[params] At most nine user-defined positional parameters can be supplied. The are bound
to the CBShell variables \$1 to \$9. The CBShell variable \$0 is bound to the name of the scriptfile.

\end{description}


\section{Commands}


\begin{description}


\item[{cbserver [{\it serveroptions}]}] :
      starts a CBserver with the specified options and connects to it

\item[connect {\it host} {\it port}] :
      connects to an already running CBserver;
      if a public CBserver is configured then it is the default host; otherwise localhost is default host;
      4001 is default port number; if the connections fails and the local operating system can start
      a CBserver, then a local CBserver is started with default settings

\item[disconnect] :
      disconnects from a CBserver

\item[stop] :
      stops the CBserver which is currently connected

\item[tell {\it frames}] :
      tells frames to the CBserver; enclose the frames in double quotes

\item[untell {\it frames}] :
      untells frames from the CBserver; enclose the frames in double quotes

\item[retell {\it untellFrames} {\it tellFrames}] :
      untells and tells frames to a CBserver in one transaction; enclose both arguments in
      double quotes\footnote{There is a small syntactic restriction for the retell
      command. You need to avoid a line consisting just of a double quote to terminate the untellFrames.
      Instead, start the tellFrames in the same line that terminates the untellFrames.}

\item[tellModel {\it file1} {\it file2} ...] :
      tells files to the CBserver; the files can have file types ".sml" and ".txt";
      if the first file from the list exists on the the computer of the CBShell client, then 
      the files will be loaded from the local file system, 
      otherwise the CBserver is requested to load the files from its
      own file system

\item[ask {\it Query} {[{\it QueryFormat} [{\it AnswerRep} [{\it RollbackTime}]]]}] :
      asks a query; possible query formats are {\tt OBJNAMES} and {\tt FRAMES};
      the answer representation can
      be {\tt LABEL}, {\tt FRAME}, {\tt FRAGMENT}, {\tt FRAGMENTswi}, {\tt JSONIC}, {\tt default}, or a user-defined answer format; the rollback
      time shall normally be set to {\tt Now}. See also subsection \ref{sec:rollbacktime} for more information.
      If the query format is {\tt OBJNAMES}, then {\it Query} is a string
      in double quotes containing a
      query call (or a comma-separated list if query calls). It the query format is {\tt FRAMES},
      then {\it Query} is a string of Telos frames including a query definition.
      The answer representations are discussed in section \ref{sec:answerformat}.

\item[hypoAsk {\it frames} {\it Query}  {[{\it QueryFormat} [{\it AnswerRep} [{\it RollbackTime}]]]}] :
      tells frames temporarily and asks a query

\item[lpicall {\it lpicall}] :
      executes the LPI call; only for debugging purposes

\item[prolog {\it prologstatement}] :
      executes the Prolog statement; only for debugging purposes

\item[why] :
      gets error messages for the last transaction and prints them on stdout
      
\item[result {\it completion} {\it result}] :
      compares the given result with the last result which has been received;
      this command hence can be used to check whether the CBserver produces
      the expected completion (ok, error) and result; use this command
      in combination with the option {\tt -l}

\item[cd {\it mod}] :
      changes the module context of this shell; if the parameter {\it mod} is omitted, 
      the module context will set the module to the user's home module, by default {\tt oHome}

\item[pwd] :
      display the current absolute module path, e.g. "System-oHome-MyModule"

\item[lm {\it mod}] :
      list all frames defined in a module;
      shortcut for {\tt ask listModule[{\it mod}/module]};
      uses currentmodule if called without parameter

\item[ls {\it class}] :
      display the instances of {\it class};
      uses {\tt Individual} as default if called without parameter

\item[mkdir {\it module}] :
      creates a new module with the given name within the current module; implemented by
      a tell operation "{\tt mod in Module end}" ; so we mimick the navigation within modules by commands known from Unix
      to manage directories

\item[showAnswer] :
      print the last result on standard output; this can be useful if you
      employ the CBserver as a generator in a shell pipe 
      (see Graphviz case below)

\item[showAnswer $>$ {\it filename}] :
      same as showAnswer but output is redirected to filename

\item[who] :
      show the list of users that have at any time been enrolled to this database;
      implemented by a query that displays the instances of the class {\tt CB\_User}

\item[sub] :
      show the list of visible submodules to which the user can branch via the 'cd' command

\item[show {\it name}] :
      show the frame with the given name; shortcut for
      {\tt ask get\_object[{\it name}/objname]}

\item[echo {\it string}] :
      prints the string to standard output; use double quotes if the string has multiple words;
      a sequence '$\backslash\backslash$n' within the string is replaced by a newline character; 

\item[echo -n {\it string}] :
      like the one-argument variant but no newline is printed after the string

\item[nl] :
      prints a newline character on standard output

\item[exit] :
      exits the shell (also stops a server which has been started in this shell)


\end{description}





Command arguments with white space characters have to be enclosed in double quotes ('"').
Command arguments may span multiple lines. Lines starting with '\#' are comment lines\footnote{
The last line of a CBShell script file should not be a comment line.
Otherwise, CBShell fails to recognize the end of file correctly.}.
If an argument contains a string of the form {\tt \$PropName}, it will be replaced with
the value of the corresponding Java property (which may be defined using the -D option
of the Java Virtual Machine), if the property is defined.
There are a couple of legacy commands that we support for backward compatibility:

\begin{description}

\item[{startServer [{\it serveroptions}]}] :
      synonym for 'cbserver'

\item[enrollMe {\it host} {\it port}] :
      synonym for 'connect'

\item[showUsers] :
      synonym for 'who'

\item[showModules] :
      synonym for 'sub'

\item[listModule {\it mod}] :
      synonym for 'lm'

\item[listClass {\it class}] :
      synonym for 'ls'

\item[getErrorMessages] :
      synonym for 'why'

\item[getModulePath] :
      synonym for 'pwd'

\item[setModule {\it mod}] :
      synonym for 'cd'

\item[stopServer] :
      synonym to 'stop'

\item[cancelMe] :
      synonym for 'disconnect'

\item[newline] :
      synonym for 'nl'

\item[quit] :
     synonym for 'exit'

\end{description}

The CBShell utility can be used in Unix pipes to extract textual output The CBserver and
pass it to subsequent programs as input. To do so, you should start the CBserver with
tracemode {\tt no} and using the {\tt showAnswer}  command of CBShell to specify the elements
to be written to standard output. The CBShell script below realizes the extraction of 
Graphviz [\url{http://graphviz.org}] specifications from ConceptBase.cc:

\begin{verbatim}
   # File: myscript
   connect
   tellModel ERD-graphviz2
   tellModel UniversityModel
   ask ShowERD[UniversityModel/erd] OBJNAMES default Now
   showAnswer
   exit
\end{verbatim}



The complete example is available from the CB-Forum at \url{http://merkur.informatik.rwth-aachen.de/pub/bscw.cgi/2519759}.
Another resource for CBShell scripts is the list of test scripts at
\url{http://merkur.informatik.rwth-aachen.de/pub/bscw.cgi/2596438}.

\subsection{Rollback time for ASK}
\label{sec:rollbacktime}

The 'ask' and 'hypoAsk' commands allow to specify a so-called rollback time for the query.
Normally the rollback time is set to 'Now', i.e.\ the current database state. You can
also specify a precise millisecond specify the database state on which you want to evaluate
the query, e.g. 'millisecond(2020,2,7,12,27,10,230)', that is 2010-02-07T12:27:10.230UTC.
Since this is a bit cumbersome, you can also specify the name of any object as 'rollbacktime'
argument. Then, the CBserver shall take the starttime of that object as the rollback time,
for example 

\begin{verbatim}
   ask get_object[AccountBill/objname] OBJNAMES FRAME transfer1
\end{verbatim}

where 'transfer1' is just another object name.




\subsection{Argument delimiters}
\label{sec:argdelim}

If an argument to a CBShell command spans multiple lines or contains several words separated by
blanks, then you have to enclose it in double quotes, being the default argument delimiter.
Double quotes are also used for {\tt String}  objects in Telos frames. These have to be
escaped like in the following example:

\begin{verbatim}
   tell "
   Peter with 
      comment about: \"This is a Telos string\"
   end
   "
\end{verbatim}


If many such frames need to be told, the escaping of Telos strings is cumbersome.
You can use the single quote as argument delimiter in such cases, making the
escaping the double quotes for the Telos string obsolete:

\begin{verbatim}
   tell '
   Peter with 
      comment about: "This is a Telos string"
   end
   '
\end{verbatim}



\section{Interactive use of CBShell}
\label{sec:cbshellinteractive}

The CBShell can be used to run a script file or in can be used in interactive mode.
In the interactive mode, the CBShell shall directly display the response to a command,
typically the answer of the ConceptBase server. Some of the shortcuts like 'why' and
'cd' are specifically defined to make the interactive mode more effective.

Another feature of the interactive mode is that queries that are represented as a single 
word can also be asked without the keyword 'ask' in front of it. CBShell will then ask
the query using 'default' as answer format, i.e.\ the ConceptBase server will decide
in which answer format to use. Below is a sample session of CBShell in interactive mode.

\begin{verbatim}
   cbshell
   This is CBShell, the command line interface to ConceptBase.cc
   [offline]>connect
   [localhost:4001]>mkdir M1
   yes
   [localhost:4001]>cd M1
   M1
   [localhost:4001]>tell "Employee in Class"
   no
   [localhost:4001]>why         
   Syntax error 1 in line 1, parser message:
   syntax error, unexpected ENDOFINPUT, expecting END or ENDMIT
   Syntax error Unable to parse Employee in Class.
   [localhost:4001]>tell "Employee in Class end"
   yes
   [localhost:4001]>tell "bill in Employee end"
   yes
   [localhost:4001]>ls Employee
   bill
   [localhost:4001]>show bill
   bill in Employee  
   end 
   [localhost:4001]>1+2
   3
   [localhost:4001]>stop
   [offline]>exit
\end{verbatim}


The term {\tt 1+2} is an example of a query that is not preceded by the 'ask' command. Note that such queries
may no contain blanks since it would split it into several words. Use quotes in such cases.
The prompt {\tt [offline]} indicates that the CBShell is not yet connected to a CBserver. Once connected,  
it displays the hostname and port number of the connected CBserver as prompt.
The CBserver is started with disabled tracing. The trace messages of the CBserver would otherwise 
be displayed in the output as well.


\section{Positional parameters}
\label{sec:cbshellvars}

A CBShell script can use variables \$0 ... \$9 inside the script to refer to 
the positional parameters supplied via the call of cbshell. Assume the following content
of the script file:

\begin{verbatim}
   # File: pascript
   cbserver -u nonpersistent -t low -port $1
   tell "$2 in Class end"
   ask find_instances[$2/class] OBJNAMES  LABEL  Now
   ask find_instances[$3/class] OBJNAMES  LABEL  Now
\end{verbatim}

and the command line call

\begin{verbatim}
   cbshell pascript 4321 MyClass Integer
\end{verbatim}


The CBShell interpreter will then replace \$1 with {\tt 4321}, \$2 with MyClass, and
\$3 with Integer.
If you supply less parameters than required by the script, it will issue an error message and quit.
If the script had started a CBserver, then this CBserver is stopped before quitting. 
Likewise, if the script had enrolled to an existing CBserver process, it will unenroll
before quitting.

The variable \$0 is bound to the name
of the script file. CBShell uses the Java tokenizer to separate the positional parameters.
The tokenizer uses white spaces (blanks, tabs) to separate tokens. Use double quotes
if one argument consists of several words, e.g. 

\begin{verbatim}
   cbshell otherscript "MyClass isA Integer end"
\end{verbatim}

An example of a script with positional parameters is provided in the CB-Forum at
\url{http://merkur.informatik.rwth-aachen.de/pub/bscw.cgi/d3364559/ticket307.cbs.txt}.
You can also call CBShell scripts within scripts/batch files of the host operating system,
see section \ref{sec:shellintegration}.


\section{Executable CBShell scripts}
\label{sec:execcbshell}

A CBShell script can be made executable under Unix/Linux and can then be used pretty much like any other
shell script. Assume that you have a CBShell script {\tt myscript} that you want to execute directly from
the command line.

As a first step, create a link to the cbshell at a common directory for installed programs:

\begin{verbatim}
   sudo ln -s $CB_HOME/cbshell /usr/bin/cbshell
\end{verbatim}

where {\tt \$CB\_HOME} is replaced by the installation directory of ConceptBase.
As a second step include the following comment as the first line of {\tt myscript}:

\begin{verbatim}
   #!/usr/bin/cbshell
\end{verbatim}

Then, make the script executable:

\begin{verbatim}
   chmod u+x myscript
\end{verbatim}

You can then directly call the script by simply typing its name:

\begin{verbatim}
   myscript
\end{verbatim}

The direct call is equivalent to the call

\begin{verbatim}
   cbshell ./myscript
\end{verbatim}
 

\section{CBShell scripts within regular shell scripts}
\label{sec:shellintegration}

CBShell offers the basic commands to interact with a ConceptBase
server. However, it lacks control structures such as loops and
conditions. It also cannot invoke arbitrary programs.
Regular shell scripts such as the Bourne shell of Unix/Linux 
do provide these capabilities, and thus it is a natural idea to
integrate CBShell scripts within regular scripts to accomplish more
sophisticated automation tasks.

As a first step, you should make the CBShell script executable and declare in its first line

\begin{verbatim}
   #!/usr/bin/cbshell -q
\end{verbatim}

This allows to pass parameters that include special characters
to the CBShell script. Assume  that the CBShell script {\tt ascript} was
declared that way. Then, it can be called in a Bourne shell like:

\begin{verbatim}
   ascript 'Jet 400' MyClass 
\end{verbatim}

The option {\tt -q}  prepares the CBShell to treat parameters with single quotes
in a special way. Assume further that {\tt ascript} has the following content:

\begin{verbatim}
   #!/usr/bin/cbshell -q
   # File: ascript
   ...
   tell "$1 in $2 end"
   ...
\end{verbatim}

CBShell will internally expand the quoted parameter {\tt 'Jet 400'} to
{\tt \textbackslash"Jet 400\textbackslash"} and then execute

\begin{verbatim}
   tell "\"Jet 400\" in MyClass end"
\end{verbatim}

The resulting frame in ConceptBase shall be

\begin{verbatim}
   "Jet 400" in MyClass end


\end{verbatim}

Essentially, single quotes of CBShell parameters are converted to double quotes within ConceptBase.
An simple example is given in
\url{http://merkur.informatik.rwth-aachen.de/pub/bscw.cgi/3372687}.
A more elaborate example is available in the CB-Forum at
\url{http://merkur.informatik.rwth-aachen.de/pub/bscw.cgi/3384265}.
The main example is in the {\tt fileSizeDemo}.


\section{CBShell in a pipe}
\label{sec:cbshellpipe}

Assume you have a program {\tt generator} that analyzes some input data (e.g from files) and
produces output in the form of CBShell commands ({\tt tell}, {\tt ask}, etc.). Then,
this output can be directly passed to CBShell in a Unix pipe:

\begin{verbatim}
   generator | cbshell -p
\end{verbatim}

The generator program must take care of generating all required commands, in particular
making sure that CBShell is connected to a CBserver. Consider as example generator the 
script file {\tt print\-files\-4cbshell}  from the
CB-Forum at
\url{http://merkur.informatik.rwth-aachen.de/pub/bscw.cgi/3384578}:

\begin{verbatim}
   #!/bin/sh
   echo "connect localhost 4001" 
   for file in *
   do
     if [ -f ${file} ] 
       then
        fsize=$( stat -c %s ${file})
        frame="'${file}' in File with size s: ${fsize} end"
        echo tell \"\\\"${file}\\\" in File with size s: ${fsize} end\"
       fi
   done
\end{verbatim}

It first generates a command to enroll to a CBserver on localhost at port 4001.
Then, it forms a frame for each filename in the current directory to tell the
file's size to the CBserver. 
You can run the script in a terminal to see the output generated by it. 

Now, assume that you have started a CBserver on localhost with port number 4001.
You should tell at least the following frame to the Cbserver to define
the class {\tt File} to which the above script refers to:

\begin{verbatim}
   File in Class with
     attribute
       size: Integer
   end
\end{verbatim}


Then execute in a terminal window:

\begin{verbatim}
   printfiles4cbshell | cbshell -p
\end{verbatim}

It will tell the file size information to the CBserver. At the end, the end of file
detection of the CBShell will trigger the CBShell to exit from the CBserver.

It is also possible to continue the pipe to a post-processing program. In this case,
the commands sent to CBShell should include {\tt ask} commands in combination
with {\tt showAnswer}. The following pipe shows the main idea:

\begin{verbatim}
   generator | cbshell -p | postprocessor
\end{verbatim}

An elaborate example of this usage is in the CB-Forum at
\url{http://merkur.informatik.rwth-aachen.de/pub/bscw.cgi/3421289}.
The example shows how to extract module import statements from program
source files, store them in ConceptBase and let ConceptBase produce
a graph specification, layed out by GraphViz.



